\documentclass[11pt]{article}
\usepackage[margin=1in]{geometry}
\usepackage[T1]{fontenc}
\usepackage[utf8]{inputenc}
\usepackage{lmodern,microtype,amsmath,amssymb,amsthm,mathtools}
\usepackage[colorlinks=true,linkcolor=blue,citecolor=blue,urlcolor=blue]{hyperref}
\hypersetup{pdftitle={Exponential parity sample complexity throughout the subcritical width regime},pdfauthor={Dmitry Vaintrob; Kaarel Hanni; OpenAI Codex}}
\setlength{\emergencystretch}{2em}
\newtheorem{theorem}{Theorem}[section]
\newtheorem{lemma}[theorem]{Lemma}
\newtheorem{proposition}[theorem]{Proposition}
\newtheorem{corollary}[theorem]{Corollary}
\theoremstyle{remark}\newtheorem{remark}[theorem]{Remark}
\DeclareMathOperator{\Var}{Var}
\DeclareMathOperator{\sign}{sign}
\DeclareMathOperator{\dist}{dist}
\DeclareMathOperator{\KL}{KL}
\newcommand{\E}{\mathbb E}
\newcommand{\Pp}{\mathbb P}
\title{Exponential parity sample complexity\\throughout the subcritical width regime}
\author{Dmitry Vaintrob \and Kaarel H\"anni \and OpenAI Codex}
\date{October 15, 2026}
\begin{document}
\maketitle
\begin{abstract}
We prove an exponential sample lower bound for learning parity with a
specified Bayesian neural-network prior throughout its subcritical scaling
regime. The network has fixed depth, Gaussian weights, $\tanh$ hidden
activations, and a linear readout. Before training, its prior is conditioned
to have a fixed absolute output margin at every Boolean input; training
then conditions on exact observed labels. For every fixed $\gamma<1$,
there is $c>0$ such that each sample size $n\le e^{ck}$ leaves at least
$1/16$ of the inputs of $k$-bit parity without correct-label posterior
probability greater than $2/3$, with probability at least $3/4$ over
independent uniform training inputs, in the width-first limit.
The proof treats the sign-conditioned Gaussian limit, the Gaussian limit
with a fixed global margin, and concentration on minimum-energy sign
functions. Its main ingredients are one-sided squared-loss sampling in
$O(D\log D)$ examples, a discrete Laplacian identity for clipped parity
correlation, sharp kernel spectra, and a conditioning estimate forcing
large output magnitudes at the fixed-margin Gaussian endpoint. The result
concerns the actual posterior confidence, includes every sample size below
the exponential threshold, and permits arbitrary weights among tied
minimum-energy functions.
\end{abstract}

\section{Model and main theorem}
Let $\mathcal X=\{-1,1\}^d$, with uniform probability measure, and let
$f(x)=\prod_{j=1}^d x_j$. We study the following precisely specified prior.
All hidden layers have width $N$, the depth $L\ge2$ is fixed, and
\[
 h^0(x)=(1,x),\qquad h^s(x)=\tanh(W^sh^{s-1}(x))\quad(1\le s<L),
 \qquad z_N(x)=W^Lh^{L-1}(x).
\]
There are no further biases. All weights are independent centered
Gaussians, with variances
\begin{equation}\label{eq:weights}
 \Var W^1_{ij}=N^{1-\gamma}/d,\qquad
 \Var W^s_{ij}=N^{-\gamma}\quad(1<s<L),\qquad
 \Var W^L_i=N^{-\gamma-1}.
\end{equation}
Fix a margin $a>0$; the convention $a=1/2$ is included. Before training,
condition on $|z_N(x)|>a$ for every $x\in\mathcal X$. Given independent
uniform inputs $S=(X_1,\ldots,X_n)$, subsequently condition on
$\sign z_N(X_i)=f(X_i)$ for all $i$. These two conditioning operations
are equivalent to conditioning on their intersection.

Define the correct-label posterior probability and its confidence-failure
fraction by
\begin{align*}
 q_{N,S}(x)&=\mathbb P\{f(x)z_N(x)>0\mid
       |z_N(y)|>a\ \forall y;\ f(X_i)z_N(X_i)>0\ \forall i\},\\
 F_N(S)&=\mathbb P_x\{q_{N,S}(x)\le2/3\}.
\end{align*}
Thus a prediction succeeds only when its actual posterior probability of
being correct is strictly greater than $2/3$. This includes equality among
failures. Expectations over a fresh input always use the uniform cube
distribution independently of the training inputs.

\begin{theorem}\label{thm:main}
For every fixed $L\ge2$ and $a>0$, there are constants $c>0$ and $d_0$
such that, for every $d\ge d_0$, every integer $0\le n\le e^{cd}$, and
every fixed $\gamma<1$,
\[
 \mathbb P_S\!\left\{\liminf_{N\to\infty}F_N(S)\ge\frac1{16}\right\}
 \ge\frac34.
\]
The conditioning probabilities are positive for all sufficiently large
$N$, at each fixed $d,L,a,\gamma,S$.
\end{theorem}

The width limit is taken at fixed dimension and fixed training sample.
The exponential statement then concerns $d=k\to\infty$. No joint limit
or quantitative minimum width is asserted here. The probability guarantee
holds for each specified sample size with uniform constants; it does not
assert a single event for every prefix of an infinite sample sequence.
In particular, the proof does not assume monotonicity of this learner's
fixed-target risk in sample size.

Parity has a Boolean circuit with $O(d)$ gates. It also has a
polynomial-size, polynomial-weight $\tanh$ realization at every fixed
depth, as shown in Appendix~\ref{sec:paritycircuit}. The theorem therefore
rules out a sample guarantee polynomial in those complexity measures for
this width-first learner. It does not claim that identifying parity from
examples is intrinsically exponentially difficult.

\subsection{Proof structure and the three limiting posteriors}
Write $\Psi(t)=2\arcsin(t)/\pi$ and $\beta=2^{-(L-1)}$. The relevant
covariance is
\[
 K(x,y)=\Psi^{\circ(L-1)}\!\left(\frac{1+x\cdot y}{d+1}\right).
\]
All population function norms use the uniform probability measure on the
cube. Write $\Pi_{\le r}g=g_{\le r}$ for the Walsh projection to degrees
at most $r$, and $g_{>r}=g-g_{\le r}$. For the positive-definite covariance
$K$, the Gaussian energy norm is $\|g\|_K^2=g^{\mathsf T}K^{-1}g$;
this uses the covariance matrix itself, without an extra factor $2^{-d}$.
At $\gamma<0$ the limit is a Gaussian process conditioned on the training
signs. At $\gamma=0$ it retains the fixed absolute margin at every input.
For $0<\gamma<1$ the actual network posterior concentrates on compatible
globally separated functions of minimum Gaussian energy; the proof needs
no formula for the weights of ties.

The common obstacle is to control low-degree functions selected after
seeing the training data. We prove a uniform squared negative-part
estimate using the stronger metric bound inside Lee's chaining argument
\cite{Lee2023}. A discrete Laplacian identity then converts an $L^2$
bound on a degree-$r$ polynomial into a clipped parity-correlation bound
with factor $r/d$. An additional short band of Fourier levels handles the
gaps between adjacent sampling thresholds. At $\gamma=0$, Gaussian
entropy bounds and an auxiliary independent sample preserve the exact
globally conditioned posterior throughout the argument.

\section{The kernel rate of the actual network}
Let $M=2^d$, and identify output functions with vectors in $\mathbb R^M$.
Define
\begin{equation}\label{eq:kernel}
 \Psi(t)=\frac2\pi\arcsin t,\qquad
 K(x,y)=\Psi^{\circ(L-1)}\!\left(\frac{1+x\cdot y}{d+1}\right).
\end{equation}
The positive definiteness of $K$ is proved in the next section.
For each Boolean function $g$, put
\begin{equation}\label{eq:energy}
 E_K(g)=\min_{g(x)z(x)\ge a\ \forall x}\frac12 z^TK^{-1}z.
\end{equation}
The minimum exists since its feasible set is closed and nonempty and
the objective is coercive. Replacing the closed inequalities by strict
ones gives the same infimum, by scaling a feasible vector by $1+\epsilon$.

\begin{proposition}[Posterior concentration on minimum-energy signs]
\label{prop:rate}
For the actual network~\eqref{eq:weights}, at every fixed
$0<\gamma<1$,
\begin{equation}\label{eq:rate}
 \lim_{N\to\infty}N^{-\gamma}\log
 \mathbb P\{g(x)z_N(x)>a\ \forall x\}=-E_K(g).
\end{equation}
Consequently, after any fixed noiseless training sample, posterior
mass tends to one on the compatible functions minimizing $E_K$.
No assertion about the relative weights of tied minimizers is needed.
\end{proposition}

\begin{proof}
For hidden layer $s$, let
$Q_N^s(x,y)=N^{-1}\sum_i h_i^s(x)h_i^s(y)$.
We first establish, for each fixed $\epsilon>0$,
\begin{equation}\label{eq:freeze}
 \mathbb P\{\|Q_N^{L-1}-K\|>\epsilon\}
       \le C_\epsilon e^{-c_\epsilon N}
\end{equation}
for all sufficiently large $N$. Any matrix norm can be used, since
$M$ is fixed. Conditional on the preceding layer, the rows of the
next hidden layer are independent bounded vectors. Hoeffding's
inequality controls all their empirical second moments. Their
preactivations have covariance $N^{1-\gamma}Q_N^{s-1}$; at the first
layer the covariance is
$N^{1-\gamma}(1+x\cdot y)/d$. Because $N^{1-\gamma}\to\infty$,
the conditional second-moment map converges to the Gaussian-sign map
\[
 Q\longmapsto \left[\Psi\!\left(
 \frac{Q(x,y)}{\sqrt{Q(x,x)Q(y,y)}}\right)\right]_{x,y}.
\]
This follows by dominated convergence for a Gaussian vector, since
each scalar variance stays bounded away from zero in a neighborhood
of the limiting covariance. The convergence is uniform in such a
neighborhood: a compactness subsequence and a continuous square root
coupling reduce the assertion to the same dominated-convergence
argument. The limiting map is continuous there. Inducting through
the fixed number of layers, with smaller fixed tolerances at earlier
layers, proves~\eqref{eq:freeze}. The first-layer limit normalizes
the diagonal $(d+1)/d$, giving the kernel in~\eqref{eq:kernel}.

Conditionally on the hidden layers, the readout is exactly Gaussian:
\begin{equation}\label{eq:readout}
 Z_N\mid Q_N^{L-1}\sim\mathcal N(0,N^{-\gamma}Q_N^{L-1}).
\end{equation}
Couple it to $Y_N=N^{-\gamma/2}K^{1/2}\xi$, where
$\xi\sim\mathcal N(0,I_M)$ is independent of the hidden layers.
On a fixed sufficiently small neighborhood of the positive-definite
matrix $K$, the square-root map is Lipschitz. Hence, for $\delta>0$,
\[
 \mathbb P\{\|Z_N-Y_N\|>\delta\}
 \le C_\epsilon e^{-c_\epsilon N}
 +\mathbb P\{\|\xi\|>\delta N^{\gamma/2}/(C\epsilon)\}.
\]
First divide logarithms by $N^\gamma$ and let $N\to\infty$;
then let $\epsilon\downarrow0$. The resulting upper limit is
$-\infty$, since $\gamma<1$. The coupled vectors are therefore
exponentially equivalent at speed $N^\gamma$.

The fixed-dimensional Gaussian vectors $Y_N$ have large-deviation
rate $I(z)=z^TK^{-1}z/2$, directly from their densities. Exponential
equivalence transfers the same bounds to $Z_N$. Applying the upper
bound to a closed margin orthant and the lower bound to its interior
gives~\eqref{eq:rate}, since their infima agree. Finally, there are
only finitely many Boolean functions for fixed $d$. Dividing each
compatible orthant probability by their sum proves the posterior
concentration claim.
\end{proof}

This is a rate statement for the network itself. It does not require
relative-error approximation of rare probabilities, nor convergence
in total variation to a Gaussian posterior.

\begin{proposition}[The nonpositive-exponent posterior limits]
\label{prop:gp-limits}
For $\gamma=0$, the Boolean posterior converges to that of
$G\sim\mathcal N(0,K)$ conditioned on $|G(x)|>a$ everywhere and on
the observed signs. For $\gamma<0$, it converges to the Gaussian
posterior conditioned only on the observed signs.
\end{proposition}
\begin{proof}
The hidden-kernel argument establishing~\eqref{eq:freeze} holds
for every $\gamma<1$. By~\eqref{eq:readout}, the normalized vector
$N^{\gamma/2}Z_N$ therefore converges in law to
$\mathcal N(0,K)$. Its all-input margin is $aN^{\gamma/2}$, which
equals $a$ when $\gamma=0$ and tends to zero when $\gamma<0$.
The limiting Gaussian is nonsingular. Thus all boundaries involved
in any fixed full-sign event or margin event have probability zero,
and the conditioning events have positive probability. Taking
ratios of their probabilities gives convergence of each Boolean
posterior mass. There are finitely many masses for fixed $d$.
\end{proof}


\section{Fourier level masses of the kernel}\label{app:spectrum}
Put $D=d+1$, $\beta=2^{-(L-1)}$, and $F=\Psi^{\circ(L-1)}$.
The characters $\chi_U(x)=\prod_{j\in U}x_j$ diagonalize the
normalized covariance operator $M^{-1}K$, $M=2^d$. Let $\lambda_j$
be its eigenvalue at degree $j$. We first work on the full $D$-cube
with kernel $F(u\cdot v/D)$, whose degree-$j$ eigenvalue is denoted
by $a_j$. Constants in this and the next section may depend on $L$.
Notation $b\asymp c$ means two inequalities by positive constant
multiples, uniformly in the stated ranges.

\begin{lemma}[Power coefficients and paired eigenvalues]\label{lem:power}
Write $F(z)=\sum_{m\text{ odd}}c_mz^m$. Every odd coefficient is
positive, $\sum_m c_m=1$, and
\[
 c_m\asymp\frac{\Gamma(m-\beta)}{m!}\asymp m^{-1-\beta}
 \qquad(m\ge1\text{ odd}).
\]
The full-cube eigenvalues vanish at even degrees and strictly decrease
through odd degrees. They are positive at every odd degree. On the
original $d$-cube,
\[
 \lambda_{2k}=\lambda_{2k+1}=a_{2k+1}>0,
\]
with a final unpaired eigenvalue if $d$ is even. In particular $K$
is positive definite.
\end{lemma}
\begin{proof}
The series of $\Psi$ has positive coefficients at every odd exponent,
with
\[
 [z^{2k+1}]\Psi(z)=\frac2\pi\frac{\binom{2k}{k}}{4^k(2k+1)}.
\]
Composition preserves strict positivity and oddness. Near $z=1$ in
a slit neighborhood,
\[
 \Psi(z)=1-b(1-z)^{1/2}+O((1-z)^{3/2}),
 \qquad b=\frac{2\sqrt2}{\pi}.
\]
Induction gives
$F(z)=1-A(1-z)^\beta+O((1-z)^{3\beta})$, with $A>0$;
at each composition its leading constant changes from $A_t$ to
$b\sqrt{A_t}$, starting from $A_1=b$.
These expansions admit algebraic coefficient transfer
\cite[Chapter VI]{flajolet2009}. Here are the analytic-domain details.
On the unit circle outside $\{-1,1\}$, positivity of every odd
coefficient gives $|\Psi(z)|<1$: equality in the triangle inequality
would require $z^2=1$. Hence the iterated functions extend
analytically through every compact arc away from these endpoints.
At $1$, restrict $\arg(1-z)$ to a closed proper subinterval of
$(-\pi,\pi)$. Each square-root composition reduces this angular
range, and its positive leading coefficient keeps its image in the
slit neighborhood for the next arcsine composition. Oddness gives
the identical argument at $-1$. These local sectors and the compact
arcs provide the required two-singularity indented domain.
The contributions of $1$ and $-1$ add at odd exponents, giving
\[
 c_m\sim\frac{2A}{-\Gamma(-\beta)}m^{-1-\beta}
 \quad(m\to\infty\text{ odd}).
\]
Together with positivity of the finitely many initial coefficients
and the gamma-ratio asymptotic, this proves the coefficient comparisons.
For $L=2$ the explicit formula and Stirling's bounds suffice instead.

Let $p_m(j)$ be the probability that the simple coordinate-flip walk
on the $D$-cube reaches a specified distance-$j$ vertex in $m$ steps.
Expanding $(u\cdot v/D)^m$ shows that its degree-$j$ Fourier coefficient
is $p_m(j)$. Thus $a_j=\sum_{m\text{ odd}}c_mp_m(j)$.
For the walk recursion,
\begin{align*}
 D[p_{m+1}(j)-p_{m+1}(j+2)]
 ={}&j[p_m(j-1)-p_m(j+1)]\\
 &+(D-j-2)[p_m(j+1)-p_m(j+3)].
\end{align*}
Induction from time zero gives nonincreasing probabilities through
each parity class; terms with zero boundary coefficients are omitted.
The term $m=j$ makes the odd-eigenvalue inequalities strict and proves
positivity. Even-degree eigenvalues vanish by parity.
Restricting the distinguished coordinate to $+1$ gives
$\lambda_q=a_q+a_{q+1}$, proving the claims.
\end{proof}

For odd $r$, the exact combined multiplicity of reduced degrees
$j-1,j$ is $\binom d{j-1}+\binom dj=\binom Dj$. Consequently
\begin{equation}\label{eq:paired}
 D_r=\sum_{\substack{j\le r\\j\text{ odd}}}\binom Dj,
 \quad
 \tau_r=\sum_{\substack{j>r\\j\text{ odd}}}\binom Dj a_j,
 \quad \lambda_{r+1}=a_{r+2}.
\end{equation}
The odd cutoff is essential for these exact identities.

\begin{lemma}[Sharp level and tail estimates]\label{lem:sharp}
Uniformly for odd $1\le j\le D/4$ and odd $1\le r\le D/8$,
\begin{align}
 \binom Dj a_j&\asymp j^{-1-\beta},\label{eq:level}\\
 \tau_r&\asymp r^{-\beta},\label{eq:trace}\\
 \frac{\lambda_{r+1}}{\tau_r}
 &\asymp\frac1{r\binom D{r+2}},\label{eq:ratio}\\
 s_{2,r}:=\sum_{q>r}\binom dq\lambda_q^2
 &\le\frac{C r^{-2\beta-2}}{\binom D{r+2}}.
 \label{eq:secondmoment}
\end{align}
The same level-mass estimate holds up to $j\le\rho D$ for any fixed
$\rho<1/2$, with constants depending on $\rho$.
\end{lemma}
\begin{proof}
The gamma integral and Lemma~\ref{lem:power} give, for odd $j$,
\begin{align}
 a_j&\asymp D^{-\beta}\int_0^\infty
 u^{-1-\beta}p(u)^j(1-p(u))^{D-j}\,du,
 \quad p(u)=\frac{1-e^{-2u}}2,\label{eq:poisson}\\
 &\asymp D^{-\beta}\int_0^{1/2}p^j(1-p)^{D-j}w_\beta(p)\,dp,
 \label{eq:betaweight}
\end{align}
where $w_\beta(p)=[-\frac12\log(1-2p)]^{-1-\beta}/(1-2p)$.
Indeed the gamma integral Poissonizes the walk time, and independent
Poisson coordinate counts give independent coordinate parities with
success probability $p(u)$. Since an odd-distance vertex has zero
probability at even walk times, no extra parity factor occurs.
The weight $w_\beta$ is integrable near $1/2$ and comparable to
$p^{-1-\beta}$ on $[0,b]$ for every fixed $b<1/2$.

Fix $\rho<b<1/2$ and take $j\le\rho D$. The integral on $[0,b]$
is comparable to the complete beta integral
$B(j-\beta,D-j+1)$. For the lower comparison, a beta random variable
with these parameters has mean at most $\rho$, so Markov's inequality
puts at least a fraction $1-\rho/b$ of its mass in $[0,b]$.
The identity
\[
 \binom Dj B(j-\beta,D-j+1)
 =\frac{\Gamma(D+1)\Gamma(j-\beta)}
        {\Gamma(j+1)\Gamma(D+1-\beta)}
 \asymp D^\beta j^{-1-\beta}
\]
gives the desired mass. On $[b,1/2]$, the binomial factor decreases
with $p$ and, after multiplication by $\binom Dj$, is at most
$\exp[-D\operatorname{KL}(\rho\Vert b)]$. The remaining integrable
weight makes this contribution exponentially small. This proves
\eqref{eq:level} and its stated extended range.

Sum \eqref{eq:poisson} over the odd distances above $r$. Dropping the
parity restriction and using the binomial Markov bound gives
\[
 \tau_r\le C D^{-\beta}\int_0^\infty
 u^{-1-\beta}\min\{1,Du/(r+1)\}\,du
 \le C(r+1)^{-\beta}.
\]
For the lower bound, sum \eqref{eq:level} over the odd $j$ in
$(r,2r+2]$, using a fixed cutoff strictly below $D/2$ as needed at
the endpoint. This yields $c r^{-\beta}$; bounded $r$ are absorbed
by a change of constants. Equation~\eqref{eq:ratio} follows from
\eqref{eq:level}, \eqref{eq:trace}, and \eqref{eq:paired}.

For the second moment, use paired multiplicities to write
$s_{2,r}=\sum_{j>r,\ j\text{ odd}}\binom Dj a_j^2$.
Up to an odd cutoff $b_D$ near $D/4$, its terms are at most
$Cj^{-2\beta-2}/\binom Dj$. Consecutive odd inverse binomial
coefficients have ratio
\[
 \frac{(j+1)(j+2)}{(D-j)(D-j-1)},
\]
uniformly less than a fixed constant below one in this range.
The sum is therefore bounded by a constant times its first term,
at $j=r+2$. Beyond $b_D$, monotonicity gives the bound
$a_{b_D}\tau_{b_D}\le C D^{-2\beta-1}/\binom D{b_D}$,
which is exponentially smaller than the claimed bound uniformly
for $r\le D/8$. This proves \eqref{eq:secondmoment}.
\end{proof}

Finally, two-step binomial ratios give, uniformly for odd $r\le D/8$,
\begin{equation}\label{eq:dimensionratios}
 \binom Dr\le D_r\le C\binom Dr,
 \qquad \frac{D_r}{\binom D{r+2}}\le C(r/D)^2.
\end{equation}
All constants used here can be fixed independently of a subsequently
chosen small positive limit $r/d\to\alpha$.

\begin{corollary}[Lower bounds used by the interpolation]\label{cor:lower-spectrum}
For sufficiently large $d$,
\[
 \lambda_0\ge c/d,\qquad
 \lambda_j\ge \frac{c d^{-1-\beta}}{\binom dj}\quad(0\le j\le .49d),
 \qquad v:=\lambda_{\min}(K)\ge c(d\log d)^{-\beta}.
\]
Also $v\to0$.
\end{corollary}
\begin{proof}
The first bound follows from the level estimate at full-cube degree one.
For odd reduced degree $j$, use $\lambda_j=a_j$ and
$\binom Dj/\binom dj=D/(D-j)$. For even $j$, use
$\lambda_j=a_{j+1}$ and
$\binom D{j+1}/\binom dj=D/(j+1)$.
The extended level bound with any fixed cutoff between $.49$ and $1/2$
then gives the stated uniform lower bound, including $j=0$.

For any odd full-cube degree $j\le D$, restrict the integral
\eqref{eq:poisson} to $u\ge\tfrac12\log(2D)$. In this range
$p(u)\ge(1-1/(2D))/2$ and $1-p(u)\ge1/2$, so
$p(u)^j(1-p(u))^{D-j}\ge c2^{-D}$. Integrating $u^{-1-\beta}$ gives
$a_j\ge c2^{-D}D^{-\beta}(\log D)^{-\beta}$.
Every reduced eigenvalue equals an odd $a_j$, and the matrix eigenvalues
are $M$ times the normalized ones. Thus $v\ge c(d\log d)^{-\beta}$.
Finally, for any fixed small positive cutoff fraction, the minimum
normalized eigenvalue is at most the average in its tail, whence
$v\le M\tau_r/(M-D_r)=O(d^{-\beta})$. This proves $v\to0$.
\end{proof}

\section{One-sided squared sampling with one logarithm}
\label{sec:sampling}

The sampling estimate needed below concerns a squared one-sided loss.
It follows from a metric estimate in the proof of Lee's chaining theorem
\cite[Lemma~2.12, equation~(2.20)]{Lee2023}. We spell out the comparison:
a first-power ReLU sampling theorem alone would not give this result.
Throughout this section, $C$ denotes a universal constant that may change
from line to line, and $t_-:=\max\{-t,0\}$.

\begin{lemma}[A squared Lipschitz process]\label{lem:squared-process}
Let $a_1,\ldots,a_m\in\mathbb R^D$, and set
\[
 R=\max_{i\le m}\|a_i\|_2,
 \qquad V=\sup_{\|u\|_2\le1}
       \left(\sum_{i=1}^m\langle a_i,u\rangle^2\right)^{1/2}.
\]
For each $i$, let $\psi_i:\mathbb R\to\mathbb R$ satisfy
$\psi_i(0)=0$ and $\operatorname{Lip}(\psi_i)\le1$.
For independent uniform signs $\epsilon_i$,
\begin{equation}\label{eq:squared-process}
 \mathbb E_\epsilon\sup_{\|u\|_2\le1}
 \left|\sum_{i=1}^m\epsilon_i\psi_i(\langle a_i,u\rangle)^2\right|
 \le C R\sqrt{\log(2(m+D))}\,V.
\end{equation}
\end{lemma}

\begin{proof}
Use the genuine seminorms $N_i(u)=|\langle a_i,u\rangle|$ and the
symmetric quasimetric
\[
 d_A(u,v)^2
 =\sum_{i=1}^m N_i(u-v)^2\bigl(N_i(u)^2+N_i(v)^2\bigr).
\]
The stronger intermediate estimate established in the proof of
\cite[Lemma~2.12]{Lee2023} is
\begin{equation}\label{eq:lee-metric}
 \gamma_2(T,d_A)
 \le C\bigl(\lambda\sqrt{\log(2D)}+\kappa\bigr)
       \sup_{u\in T}\left(\sum_i N_i(u)^2\right)^{1/2},
 \qquad T\subseteq B_2^D,
\end{equation}
where
\[
 \lambda=\max_i\bigl(\mathbb E N_i(g)^2\bigr)^{1/2},
 \qquad \kappa=\mathbb E\max_i N_i(g),
 \qquad g\sim\mathcal N(0,I_D).
\]
Here $\gamma_2$ is the generic-chaining functional. This is equation
(2.20) of that proof, with the energy normalization (2.17) undone.
Seminorms follow from the norm statement by replacing
$N_i(u)$ with $(N_i(u)^2+\varepsilon^2\|u\|_2^2)^{1/2}$ and then
letting $\varepsilon\downarrow0$; the resulting metrics dominate
$d_A$, and the parameters on the right converge.
For our rank-one seminorms,
$\lambda=R$ and $\kappa\le R\sqrt{2\log(2m)}$.

Let $X_u=\sum_i\epsilon_i\psi_i(\langle a_i,u\rangle)^2$.
Its subgaussian increment metric is bounded by $\sqrt2\,d_A$, because
\begin{align*}
 &\sum_i\left(\psi_i(\langle a_i,u\rangle)^2
             -\psi_i(\langle a_i,v\rangle)^2\right)^2\\
 &\quad\le\sum_i |\langle a_i,u-v\rangle|^2
       \bigl(|\langle a_i,u\rangle|+|\langle a_i,v\rangle|\bigr)^2
 \le 2d_A(u,v)^2.
\end{align*}
We used both Lipschitz continuity and
$|\psi_i(t)|\le|t|$. The subgaussian chaining inequality, together
with $X_0=0$, bounds $\mathbb E\sup X_u$ by
$C\gamma_2(B_2^D,d_A)$. Apply the same argument to $-X_u$ and add.
Equation~\eqref{eq:lee-metric} now proves
\eqref{eq:squared-process}.
\end{proof}

\begin{remark}[The normalization in the metric estimate]
The proof of \cite[Lemma~2.12]{Lee2023} passes through a coarse additive
constant in its equation (2.26), after normalizing
$\sup_T\sum_iN_i^2=1$. The homogeneous version
\eqref{eq:lee-metric} follows by retaining the scale in the two geometric
terms of that argument. Indeed, in the rank-one case $T\subseteq B_2^D$
implies $N_i(u)\le\lambda$ on $T$; writing
$\|u\|_{\mathcal N}=\max_iN_i(u)$ gives
$d_A(u,v)\le\sqrt2\|u-v\|_{\mathcal N}$ on $T$.
Thus $\operatorname{diam}_{d_A}(T)\le C\lambda$ and the volumetric
packing constant in its Corollary~2.11 can be taken to be
$C\lambda$, since $T\subseteq\lambda B_{\mathcal N}$.
The additive term is therefore $C\lambda$, absorbed by
$C(\lambda\sqrt{\log(2D)}+\kappa)$.
The growth-functional portion of the proof is unchanged.
If the packing constant is less than one, first rescale the quasimetric
to apply that corollary in its stated normalization, and then scale back.
Most importantly, throughout that argument the $N_i$ are symmetric
seminorms. No assertion that a positive part is a seminorm is used.
\end{remark}

\begin{lemma}[Uniform one-sided squared sampling]
\label{lem:squared-sampling}
Let $(Z,Y)$ be a random pair with $Z\in\mathbb R^D$, $Y\in\{-1,1\}$,
$\mathbb E ZZ^{\mathsf T}=I_D$, and $\|Z\|_2^2=D$ almost surely.
For independent copies $(Z_i,Y_i)$, let
\[
 \Delta_n=\sup_{\|u\|_2\le1}
 \left|\frac1n\sum_{i=1}^n(Y_i\langle Z_i,u\rangle)_-^2
          -\mathbb E(Y\langle Z,u\rangle)_-^2\right|,
 \qquad h_n=\sqrt{\frac{D\log(2(n+D))}{n}}.
\]
Then $\mathbb E\Delta_n\le C h_n(1+C h_n)$.
In particular, for any fixed $\varepsilon,\delta>0$ there is a finite
$A_{\varepsilon,\delta}$ such that
\[
 n\ge A_{\varepsilon,\delta}D\log(2D)
 \quad\Longrightarrow\quad
 \mathbb P\{\Delta_n>\varepsilon\}\le\delta.
\]
The labels may depend arbitrarily on the features.
\end{lemma}

\begin{proof}
Put $Q_n=n^{-1}\sum_iZ_iZ_i^{\mathsf T}$.
Conditional on the observations, apply
Lemma~\ref{lem:squared-process} to
$a_i=Y_iZ_i/\sqrt n$ and $\psi_i(t)=t_-$.
Here $R=\sqrt{D/n}$ and $V=\sqrt{\|Q_n\|_{\mathrm{op}}}$.
Symmetrization gives
\begin{equation}\label{eq:sampling-symmetrization}
 \mathbb E\Delta_n
 \le C h_n\,\mathbb E\sqrt{\|Q_n\|_{\mathrm{op}}}.
\end{equation}
We also apply the same process lemma with $\psi_i(t)=t$.
Since $\mathbb E Q_n=I_D$, symmetrization and Jensen's inequality imply,
with $A=\mathbb E\|Q_n\|_{\mathrm{op}}$,
\[
 A\le1+\mathbb E\|Q_n-I_D\|_{\mathrm{op}}
 \le1+C h_n\mathbb E\sqrt{\|Q_n\|_{\mathrm{op}}}
 \le1+C h_n\sqrt A.
\]
Solving this quadratic inequality yields $\sqrt A\le1+C h_n$,
after increasing $C$. Substitute into
\eqref{eq:sampling-symmetrization}.
Finally, at $n\ge A_{\varepsilon,\delta}D\log(2D)$ the quantity
$h_n$ can be made uniformly small by increasing
$A_{\varepsilon,\delta}$. Markov's inequality proves the probability
statement. It is a statement for each specified sample size $n$.
\end{proof}

\begin{lemma}[Parity balances squared amplitudes]
\label{lem:parity-squared-balance}
Let $f(x)=\prod_{j=1}^d x_j$ on the uniform discrete cube, and let $P$
have Walsh degree at most $r$. Then
\begin{equation}\label{eq:parity-squared-balance}
 \mathbb E(fP)_-^2
 \ge\left(\frac12-\sqrt{\frac rd}\right)\|P\|_2^2.
\end{equation}
Consequently, for each fixed $\delta>0$ there is a universal
$A_\delta$ such that, if $r\le d/64$ and
$n\ge A_\delta D_r\log(2D_r)$, where
$D_r=\sum_{j=0}^r\binom dj$, then, with probability at least $1-\delta$,
\begin{equation}\label{eq:wrong-square-coercivity}
 \frac1n\sum_{i=1}^n(f(X_i)P(X_i))_-^2
 \ge\frac14\|P\|_2^2
 \quad\text{for every polynomial $P$ of degree at most $r$}.
\end{equation}
\end{lemma}

\begin{proof}
Write $x^{(j)}$ for $x$ with coordinate $j$ flipped, and
$D_jP=(P(x)-P(x^{(j)}))/2$. Parity changes sign under each flip.
For $F(t)=t|t|$, the inequality
$|F(a)-F(b)|\le(|a|+|b|)|a-b|$ and Cauchy--Schwarz give
\[
 \bigl|\mathbb E[fP|P|]\bigr|
 =\frac12\left|\mathbb E f(x)
      \bigl(F(P(x))-F(P(x^{(j)}))\bigr)\right|
 \le2\|P\|_2\|D_jP\|_2.
\]
Average over $j$ and use
$\sum_j\|D_jP\|_2^2\le r\|P\|_2^2$ to obtain
$\bigl|\mathbb E[fP|P|]\bigr|\le2\sqrt{r/d}\,\|P\|_2^2$.
Since $(fP)_-^2=(P^2-fP|P|)/2$, this proves
\eqref{eq:parity-squared-balance}.

The feature vector consisting of all Walsh characters of degree at
most $r$ is isotropic and has squared norm exactly $D_r$.
When $r\le d/64$, the population lower bound is at least
$3\|P\|_2^2/8$. Apply Lemma~\ref{lem:squared-sampling} with
$Y=f(X)$ and $\varepsilon=1/8$, and use homogeneity, to obtain
\eqref{eq:wrong-square-coercivity}.
\end{proof}

\section{Random tail covariance matrices}\label{app:gram}
Let $Q$ be the covariance kernel retaining only degrees greater than an
odd cutoff $r$, and write $\tau=\tau_r$, $\ell=\lambda_{r+1}$, and
$s_2=s_{2,r}$. This section proves the sample event needed later, without
assuming that a Gaussian approximation remains valid after conditioning.
The entire argument concerns the exact limiting Gaussian covariance.

\begin{lemma}[Far-pair decay]\label{lem:far}
There are constants $c_0,C>0$ such that, when $r\le D/8$ is odd,
\[
 \max_{d/4\le\dist(x,y)\le3d/4}|Q(x,y)|
 \le\eta_r:=C\sqrt D\,e^{-r/8}+Ce^{-c_0D}.
\]
\end{lemma}
\begin{proof}
The full-cube Krawtchouk polynomial satisfies
\[
 \sum_{j=0}^D\mathsf K_j(h;D)z^j=(1+z)^{D-h}(1-z)^h.
\]
If $\delta D\le h\le(1-\delta)D$, put $b=\min(h,D-h)$.
On $|z|=t$ the generating function is bounded in absolute value by
\[
 (1+t^2)^b(1+t)^{D-2b}
 \le(1+t)^D\exp\!\left[-\frac{2\delta Dt}{(1+t)^2}\right].
\]
Cauchy's coefficient bound at $t=j/(D-j)$ and the uniform Stirling
lower bound for $\binom Dj$ therefore give, for $1\le j<D$,
\begin{equation}\label{eq:kraw}
 \frac{|\mathsf K_j(h;D)|}{\binom Dj}
 \le C\sqrt D\exp[-2\delta j(1-j/D)].
\end{equation}
For reduced-cube points at distance in $[d/4,3d/4]$, their augmented
full-cube distance lies in $[D/8,7D/8]$ for large $d$.
Because $r$ is odd, the exact tail is
$Q(x,y)=\sum_{j>r,\ j\text{ odd}}a_j\mathsf K_j(h;D)$.
For $r<j\le D-r$, \eqref{eq:kraw} gives a bound
$C\sqrt D e^{-r/8}$ after summation, since the total spectral mass
is one. For $j>D-r$, monotonicity and an odd cutoff $b_D$ near $D/4$
give
\[
 \sum_{j>D-r}\binom Dj a_j
 \le a_{b_D}\sum_{k<r}\binom Dk
 \le C D^{-\beta-1/2}
       \exp[-D(h(1/4)-h(1/8))+O(1)].
\]
Here $h(t)=-t\log t-(1-t)\log(1-t)$ denotes the binary entropy in nats.
The positive entropy gap absorbs the polynomial factor and proves the
claim. Adjust constants for bounded $d$.
\end{proof}

\begin{lemma}[A tail Gram deviation bound]\label{lem:gram-general}
Let $X_1,\ldots,X_n$ be independent uniform cube points, with
$1\ll n\le M/4$. Let
\[
 B_{\mathrm{near}}=\#\{y:\dist(x,y)\notin[d/4,3d/4]\}
 \le 2e^{d h(1/4)},\qquad q=2\log M.
\]
The count is independent of $x$. For every $z>0$,
\begin{align}
 &\mathbb P\{\|Q_{SS}-\tau I_n\|_{\mathrm{op}}>z\tau\}
 \notag\\
 &\quad\le\frac C{z^2}
 \left[\frac{q\eta_r}{\tau}
       +\frac{\sqrt{nq s_2}}{\tau}
       +\frac{n\ell}{\tau}
       +\frac{n(1+B_{\mathrm{near}})}M\right]^2
       +\frac{n^2(1+B_{\mathrm{near}})}M.
 \label{eq:gram-bound}
\end{align}
Here the sampled matrix includes repeated training locations.
\end{lemma}
\begin{proof}
Let $J=Q-\tau I_M$, a zero-diagonal Hermitian matrix. Form a deterministic
mask $J_{\mathrm{far}}$ retaining only pairs at distances in
$[d/4,3d/4]$. Since $Q$ is positive semidefinite with constant diagonal
$\tau$, every entry has absolute value at most $\tau$. Thus
\begin{align*}
 \|J_{\mathrm{far}}\|_{\max}&\le\eta_r,\\
 \|J_{\mathrm{far}}\|_{1\to2}^2&\le M s_2,\\
 \|J_{\mathrm{far}}\|_{\mathrm{op}}
 &\le M\ell+\tau+\tau B_{\mathrm{near}}.
\end{align*}
For the second inequality, Parseval gives squared column norm
$M s_2-\tau^2$ for $J$, and entrywise masking only decreases it.
For the third, $\|J\|\le M\ell+\tau$, while the removed matrix has
maximum absolute row sum at most $\tau B_{\mathrm{near}}$.
The entropy bound on $B_{\mathrm{near}}$ is the elementary binomial
tail bound on distances to a fixed cube point.

Let $R$ retain each of the $M$ coordinates independently with
probability $p=2n/M$. The zero-diagonal case of Tropp's random
principal-submatrix theorem \cite[Theorem 1.1]{tropp2008}, at moment
order two and $q=2\log M$, states
\[
 (\mathbb E\|RJR\|^2)^{1/2}
 \le C\left[q(\mathbb E\|RJR\|_{\max}^2)^{1/2}
       +\sqrt{pq}(\mathbb E\|JR\|_{1\to2}^2)^{1/2}
       +p\|J\|\right].
\]
Apply it to $J_{\mathrm{far}}$ and substitute the preceding bounds:
\begin{equation}\label{eq:bernoulli-gram}
 (\mathbb E\|R J_{\mathrm{far}}R\|^2)^{1/2}
 \le C[q\eta_r+\sqrt{nq s_2}+n\ell
               +(n/M)\tau(1+B_{\mathrm{near}})].
\end{equation}
The logarithm here is that of the original dimension $M$, not the
sample size.

To pass to a uniform subset of size $n$, use a uniform permutation of
all cube points, independently of $B\sim\operatorname{Bin}(M,2n/M)$.
Its first $B$ elements form the Bernoulli subset. On $B\ge n$, the
first $n$ principal submatrix is a compression of the larger one, so
its operator norm cannot be larger. The size-$n$ event is independent
of $B$, and $\mathbb P(B\ge n)\ge1/2$ for large $n$. Thus Markov's
inequality applied to \eqref{eq:bernoulli-gram} gives the same fixed-size
probability bound with a constant factor loss.

Conditional on no collisions, an independent size-$n$ sample is a
uniform size-$n$ subset, up to an irrelevant ordering. Collisions have
probability at most $n(n-1)/(2M)$. The probability of any sampled pair
violating the distance condition is at most $n^2B_{\mathrm{near}}/M$.
On their complement, the sampled matrices $J$ and $J_{\mathrm{far}}$
agree. Adding these exceptional probabilities proves
\eqref{eq:gram-bound}.
\end{proof}


\section{Clipped parity and adaptive spectral cutoffs}
\label{sec:adaptive}
All function norms in this section use the uniform cube probability
measure. Write $\operatorname{clip}(t)=\max\{-1,\min\{1,t\}\}$ and
$\Delta=\sum_{j=1}^d(I-\mathrm{flip}_j)/2$. Thus $\Delta$ has eigenvalue
$j$ on Fourier level $j$.

\begin{lemma}[Clipped parity correlation]\label{lem:clipped-parity}
For any real cube function $P$ and any $b>0$,
\begin{equation}\label{eq:laplacian-clip}
 \left|\mathbb E f\operatorname{clip}(P/b)\right|
 \le\frac{\|\Delta P\|_1}{bd}.
\end{equation}
In particular, if $P$ has degree at most $u$, the right side is at most
$u\|P\|_2/(bd)$. If moreover $P$ has degree at most $u$, $z=P+t$, and
$|z|\ge b$ everywhere, then
\begin{equation}\label{eq:margin-clip}
 \left|\mathbb E f\operatorname{sign}z\right|
 \le\frac{u\|P\|_2}{bd}+\frac{\|t\|_2}{b}.
\end{equation}
\end{lemma}
\begin{proof}
Let $R$ flip one uniformly chosen coordinate, put $q=fP$, and set
$e=2f\Delta P/d$. Since $f$ changes sign along every edge,
$Rq=-q+e$. Put $A=\mathbb E|q+b|$ and $B=\mathbb E|q-b|$.
Jensen's inequality and invariance of the uniform measure give
$\mathbb E|Rq+b|\le A$, whereas the triangle inequality gives
$\mathbb E|Rq+b|\ge B-\|e\|_1$. Replacing $b$ by $-b$ gives the
opposite inequality, so $|A-B|\le\|e\|_1$. The identity
\[
 \operatorname{clip}(q/b)=\frac{|q+b|-|q-b|}{2b}
\]
proves \eqref{eq:laplacian-clip}. For a degree-$u$ function,
$\|\Delta P\|_1\le\|\Delta P\|_2\le u\|P\|_2$.
Finally, $\operatorname{sign}z=\operatorname{clip}(z/b)$ under the
margin assumption. Clipping is $1$-Lipschitz, so its difference from
$\operatorname{clip}(P/b)$ has mean absolute value at most
$\|t\|_1/b\le\|t\|_2/b$.
\end{proof}

The following elementary polynomial-threshold obstruction is due to
Aspnes, Beigel, Furst, and Rudich~\cite{ABFR}; we include its proof.
\begin{lemma}[Polynomial thresholds cannot approximate parity]
\label{lem:ptf-parity}
For $0\le s<d$, every degree-$s$ polynomial threshold, with a fixed common convention
for zero ties, has parity error at least
\begin{equation}\label{eq:ptf-parity}
 2^{-d}\sum_{j=0}^{\lfloor(d-s-1)/2\rfloor}\binom dj.
\end{equation}
This lower bound tends to $1/2$ whenever $s=o(\sqrt d)$.
\end{lemma}
\begin{proof}
A sufficiently small constant perturbation resolves any zero ties
without changing the specified threshold classifier. Thus suppose
$P$ has no zeros. Put $t=\lfloor(d-s-1)/2\rfloor$. If its error set
had fewer than $\sum_{j\le t}\binom dj$ points, a nonzero polynomial
$Q$ of degree at most $t$ could be chosen to vanish on every error.
The function $fPQ^2$ would then be nonnegative everywhere and positive
somewhere. But $PQ^2$, after multilinear reduction, has degree at most
$s+2t<d$, so its parity Fourier coefficient
$\mathbb E fPQ^2$ is zero, a contradiction. This proves
\eqref{eq:ptf-parity}. The last assertion follows from the central
binomial bound $\max_j2^{-d}\binom dj=O(d^{-1/2})$.
\end{proof}

\begin{lemma}[The small-sample cutoff]\label{lem:small-cutoff}
Let $s$ be an odd integer with $s\sim d^{1/3}$. For every fixed
$C,m>0$, uniformly over $0\le n\le\exp[(\log(d+1))^3]$,
\[
 C d^m(1+n)\lambda_{s+1}=o(1).
\]
Consequently, a globally margin-$a$ function $z$ satisfying
$\|z\|_K^2\le C d^m(1+n)$ has parity error at least $1/2-o(1)$,
uniformly over such functions and sample sizes. Here
$\|z\|_K^2=z^{\mathsf T}K^{-1}z$.
\end{lemma}
\begin{proof}
Lemma~\ref{lem:sharp} gives
$\lambda_{s+1}\le C s^{-1-\beta}/\binom{d+1}{s+2}$.
The logarithm of this binomial coefficient is of order
$d^{1/3}\log d$, which dominates $(\log(d+1))^3+O(\log d)$.
For $P=z_{\le s}$, spectral orthogonality therefore gives
$\|z-P\|_2^2\le\lambda_{s+1}\|z\|_K^2=o(1)$.
Because $|z|\ge a$, the signs of $z$ and $P$ can disagree on at
most $\|z-P\|_2^2/a^2=o(1)$ of the cube. Apply
Lemma~\ref{lem:ptf-parity} with $s=o(\sqrt d)$.
\end{proof}

For an odd integer $r$, retain the notation $D_r,\tau_r,s_{2,r}$
from \eqref{eq:paired} and Lemma~\ref{lem:sharp}, and write
$\ell_r=\lambda_{r+1}$.

\begin{lemma}[Adaptive cutoff and a short additional band]
\label{lem:adaptive-cutoff}
Fix $A\ge1$, $\delta>0$, and $0<\alpha<1/256$.
There exists $c>0$ such that the following holds uniformly for
\[
 \exp[(\log(d+1))^3]\le n\le e^{cd}
\]
as $d\to\infty$. Let $r$ be the largest odd integer satisfying
$A D_r\log(2D_r)\le n$, and put
\[
 \rho=\frac r{d+1},\qquad H_r=\log(e/\rho),\qquad
 u=r+2\lceil10\log(d+1)\rceil.
\]
Then
\begin{align}
 &\tfrac14(\log d)^2\le r\le\alpha d,\quad
 u/r=1+o(1),\quad u\le d/64,\label{eq:adaptive-range}\\
 &n<A D_{r+2}\log(2D_{r+2}),\qquad
 \frac{n\ell_r}{\tau_r}\le C_L A H_r,\label{eq:adaptive-ell}\\
 &\frac{nd\,s_{2,r}}{\tau_r^2}\le C_L A\rho^{-1}H_r,
 \qquad \frac{\tau_u}{\tau_r}=1+o(1),
 \qquad n\ell_u\le d^{-10}\tau_r.\label{eq:adaptive-tail}
\end{align}
Moreover, with probability at least $1-\delta$,
\begin{equation}\label{eq:adaptive-gram}
 \|(K_{>r})_{SS}\|_{\mathrm{op}}\le b_r\tau_r,
 \qquad b_r=C_{L,\delta}
     \left[1+\sqrt{A\rho^{-1}H_r}+A H_r\right].
\end{equation}
The estimates and the probability are uniform over each specified
$n$ in the range; no simultaneous event over all $n$ is asserted.
\end{lemma}
\begin{proof}
The bound $D_r\le (ed/r)^r$ for $r\le d/2$, with an immaterial
constant enlargement at bounded $r$, shows that the sampling threshold
is met by an odd integer just above $(\log d)^2/4$.
Conversely, binomial entropy bounds show that $D_r>e^{cd}$ whenever
$r>\alpha d$, if $c$ is sufficiently small. Maximality of $r$ then
gives the next-threshold inequality in \eqref{eq:adaptive-ell}.
Using \eqref{eq:dimensionratios} and
$\log(2D_{r+2})\le C r\log(e(d+1)/r)$ in that inequality yields
\[
 n\le C A\binom{d+1}{r+2}rH_r.
\]
Equations \eqref{eq:ratio} and \eqref{eq:secondmoment} now prove
the bounds involving $\ell_r$ and $s_{2,r}$.

For the additional band, apply the two-sided level estimate only at
its endpoints, to obtain
\[
 \frac{\ell_u}{\ell_r}
 \le C_L\left(\frac{r+2}{u+2}\right)^{1+\beta}
      \frac{\binom{d+1}{r+2}}{\binom{d+1}{u+2}}
 \le C_L4^{-(u-r)/2}.
\]
The last inequality follows by multiplying the exact two-step
binomial ratios below $(d+1)/8$. In particular the fixed comparison
constant $C_L$ is used once, not once per level.
Together with \eqref{eq:adaptive-ell} and the choice of $u$, this
implies $n\ell_u\le d^{-10}\tau_r$ for all sufficiently large $d$.
Also
$0\le\tau_r-\tau_u\le C_L(u-r)r^{-1-\beta}$,
so \eqref{eq:trace} gives $\tau_u/\tau_r=1+o(1)$ uniformly.

Finally apply Lemma~\ref{lem:gram-general}. Its far-coherence term
$q\eta_r/\tau_r$ tends to zero uniformly by
Lemma~\ref{lem:far} and the lower bound on $r$.
Choose $c$ also to satisfy
$2c<\log2-h(1/4)$, so that its near-pair and collision terms vanish.
The two remaining terms are bounded by
$C_L[\sqrt{A\rho^{-1}H_r}+A H_r]$.
Increasing $C_{L,\delta}$ in \eqref{eq:adaptive-gram} makes the
failure probability at most $\delta$ for all sufficiently large $d$.
\end{proof}

\section{Conditioning cost and amplitude at the Gaussian endpoint}
\label{sec:zero-conditioning}
Fix $a>0$. For an input sample $S$ and arbitrary labels $y_i\in\{-1,1\}$,
define the full event
\[
 A_S=\{|G(x)|>a\ \forall x\in\mathcal X;\ y_iG(X_i)>0\ \forall i\},
 \qquad \mathcal D(S)=-\log\mathbb P\{G\in A_S\},\quad G\sim N(0,K).
\]
All probabilities below concern this exact event. Write $M=2^d$ and
$v=\lambda_{\min}(K)>0$ for the smallest \emph{matrix} eigenvalue;
it differs from the normalized Walsh eigenvalues $\lambda_j$.

\begin{lemma}[Adaptive conditioning cost]\label{lem:zero-denominator}
For $1\le n\le e^{d/25}$, with probability at least $1-e^{-c d}$ over
independent uniform inputs, simultaneously for all labelings and every
odd $1\le r\le d/200$,
\begin{align}
 \mathcal D(S)&\le C\{d^2+n d^{1+\beta}\tau_r+D_r\log(r+1)\},
 \label{eq:zero-adaptive-cost}\\
 \mathcal D(S)&\le C\{d^2+n d^{1+\beta}\}.
 \label{eq:zero-rough-cost}
\end{align}
The constants depend only on $L,a$.
\end{lemma}
For a covariance $J$ on the cube, use the Gaussian Cameron--Martin norm $\|h\|_J^2=h^TJ^{-1}h$. For a translation-invariant kernel it equals $\sum_U\widehat h(U)^2/\lambda^J_{|U|}$, using normalized Walsh coefficients.

To prove the lemma, let $Q$ retain degrees greater than $r$, let $K_{\rm low}=K-Q$, and set $\tau=\tau_r$. Define
\[
 \widetilde K=\tau K_{\rm low}+Q.
\]
Lemma~\ref{lem:sharp} proves $\tau\asymp r^{-\beta}$, while Corollary~\ref{cor:lower-spectrum} gives $\lambda_0\ge c/d$ and $v\ge c(d\log d)^{-\beta}$. In particular $0<\tau<1$ for large $d$. The diagonal of $\widetilde K$ is $\tau(1-\tau)+\tau\le2\tau$. Its low-degree eigenvalues are those of $K$ multiplied by $\tau$, while all remaining eigenvalues agree. The Gaussian densities thus satisfy, pointwise,
\begin{equation}\label{eq:zero-density-compare}
 \frac{dN(0,K)}{dN(0,\widetilde K)}(g)
 =\tau^{D_r/2}\exp\!\left[\frac12g^T(\widetilde K^{-1}-K^{-1})g\right]
 \ge\tau^{D_r/2}.
\end{equation}
The comparison covariance is only a proof device; it does not replace the prior in the theorem.

\begin{lemma}[Localized interpolation]\label{lem:zero-bump}
For all sufficiently large $d$, with probability at least $1-e^{-cd}$ over an independent uniform sample with $1\le n\le e^{d/25}$, simultaneously for every labeling, every odd $1\le r\le d/200$, and every $A>0$, there is a function $h$ such that
\[
 |h(x)|>A\quad\forall x,\qquad y_i h(X_i)>A\quad\forall i,
 \qquad \|h\|_{\widetilde K}^2\le C A^2(d/\tau+nd^\beta).
\]
\end{lemma}
\begin{proof}
For $z\in\mathcal X$ define
\[
 b_z(x)=20\int_{2/5}^{9/20}(1-2p)^{\dist(x,z)}\,dp.
\]
Then $b_z(z)=1$, $0\le b_z(x)\le(1/5)^{\dist(x,z)}$, and its Walsh coefficients have magnitudes
\[
 \widehat b_j=20\int_{2/5}^{9/20}(1-p)^{d-j}p^j\,dp,
 \qquad \binom dj\widehat b_j\le\frac{20}{d+1},\qquad \widehat b_j\le(3/5)^d.
\]
The middle inequality follows by extending to the complete beta integral. Corollary~\ref{cor:lower-spectrum} gives
\[
 \lambda_j\ge\frac{c d^{-1-\beta}}{\binom dj}\quad(j\le .49d).
\]
Therefore the contribution of these degrees to $\|b_z\|_K^2$ is at most
\[
 C d^{1+\beta}\sum_{j\le .49d}\left[\binom dj\widehat b_j\right]^2\le C d^\beta.
\]
For $j>.49d$ the integrand increases with $p$ on the integration interval, giving $\widehat b_j\le .55^{d-j}.45^j$. Since every $\lambda_j\ge v/M$, the remaining energy is at most
\[
 \frac{2^d}{v}\sum_{j>.49d}\binom dj .55^{2(d-j)}.45^{2j}\le \operatorname{poly}(d)e^{-c_*d},
\]
where
\[
 c_*=-[\log2+h(.49)+1.02\log(.55)+.98\log(.45)]>.0062.
\]
Here $h(t)=-t\log t-(1-t)\log(1-t)$, and the binomial entropy bound applies because the exponent decreases for $j/d\ge .49$.

The low-degree energy increases by at most
\[
 \tau^{-1}D_r( .72)^d/v
\]
under damping. Uniformly in the stated range, $D_r\le\exp[(h(1/200)+o(1))d]$. Since $h(1/200)<.04$, this term decays exponentially. Hence $A_d:=\|b_z\|_{\widetilde K}^2\le Cd^\beta$.

Let $J(z,w)=\langle b_z,b_w\rangle_{\widetilde K}$. Its normalized Fourier eigenvalues are $\nu_j=\widehat b_j^2/\lambda^{\widetilde K}_j\le .72^d/(\tau v)$. Parseval and Cauchy--Schwarz give
\[
 \E_{X,Y}|J(X,Y)|\le\left(\sum_j\binom dj\nu_j^2\right)^{1/2}
 \le\sqrt{\nu_{\max} A_d}.
\]
For every subset $T$ of the sampled indices, even one chosen after seeing the sample,
\[
 \left\|\sum_{i\in T}b_{X_i}\right\|_{\widetilde K}^2
 \le nA_d+\sum_{i\ne j}|J(X_i,X_j)|.
\]
The expected cross sum is at most $n(n-1)\sqrt{\nu_{\max}A_d}$. Markov's inequality at threshold $nd^\beta$ has exponentially small failure when $n\le e^{.04d}$, because $\tfrac12\log(1/.72)>.16$. Thus all subsets simultaneously have energy at most $Cnd^\beta$. A union bound over the at most $d$ cutoffs preserves exponentially small failure; the estimates are uniform because $\tau_r\ge c d^{-\beta}$.

Also, with exponentially high probability, every pair of different sample indices has inputs at distance greater than $.2d$: a union bound costs at most $n^2\exp[-d(\log2-h(.2))]$, with $\log2-h(.2)>.19$. This event excludes repetitions. At any input at most one center is within distance $.1d$, and the contribution of all other centers is at most $n(.2)^{.1d}=o(1)$. Thus the bump sum over any selected centers is $1+o(1)$ at a selected center and at most $.2+o(1)$ everywhere else.

For $T=\{i:y_i=-1\}$ take $h=2A(1-2\sum_{i\in T}b_{X_i})$. The preceding bounds give its strict global and training margins. Since $\|1\|_{\widetilde K}^2=1/(\tau\lambda_0)\le Cd/\tau$, the squared triangle inequality gives the asserted energy bound.
\end{proof}

\begin{corollary}[A strict-margin comparator]\label{lem:comparator}
On the same sample event, for every labeling and every $a_0>0$ there
is a function $h$ with $|h(x)|>a_0$ at every input,
$y_i h(X_i)>a_0$ at every training location, and
\[
 \|h\|_K^2\le C a_0^2(d+nd^\beta).
\]
\end{corollary}
\begin{proof}
Apply Lemma~\ref{lem:zero-bump} at the fixed cutoff $r=1$, for which
$\tau_1\asymp1$. Since $\widetilde K\preceq K$, its Gaussian energy
bounds the original energy from above:
$\|h\|_K^2\le\|h\|_{\widetilde K}^2$. The same event therefore also
supplies the undamped comparator used below.
\end{proof}

\begin{proof}[Proof of Lemma~\ref{lem:zero-denominator}]
Put $t_0=\sqrt{4\tau\log(4M)}$ and $A=a+t_0+1$. Under $N(0,\widetilde K)$ the symmetric event $B=\{\max_x|G(x)|\le t_0\}$ has probability at least $1/2$, by the diagonal bound and a Gaussian union bound. Lemma~\ref{lem:zero-bump} makes $h+B$ a subset of $A_S$. Because $a$ is fixed, $\tau\ge c d^{-\beta}$, and $\beta<1$, we have $A^2\le C_a\tau d$ for large $d$. Hence
\[
 \|h\|_{\widetilde K}^2\le C(d^2+n d^{1+\beta}\tau).
\]
For a symmetric event $B$, the Gaussian translation formula, symmetry, and $\cosh t\ge1$ imply
\[
 \Pp_{\widetilde K}(h+B)\ge \Pp_{\widetilde K}(B)e^{-\|h\|_{\widetilde K}^2/2}.
\]
Combine this with \eqref{eq:zero-density-compare} and
$\log(1/\tau)\le C\log(r+1)$ to obtain
\eqref{eq:zero-adaptive-cost}. The identical argument without damping
uses $K$ itself, a buffer of squared size $O(d)$, constant-function
energy $O(d)$, and bump energy $O(nd^\beta)$, proving
\eqref{eq:zero-rough-cost}.
\end{proof}


\subsection{Large magnitudes for actual output draws}
Write $\Phi$ for the standard normal distribution function and
$\overline\Phi=1-\Phi$.
\begin{lemma}[A finite Gaussian amplitude estimate]\label{lem:zero-auxiliary}
Let $G$ be a nondegenerate centered Gaussian vector on a finite set. Suppose an event $A$ implies $|G(x)|>a$ at every coordinate and $-\log\Pp(A)\le D$, with $D\ge1$. Let $T=(x_1,\ldots,x_m)$ be distinct coordinates with $K_T\succeq wI_m$, where $w>0$. Assume $D<m$ and $e^{-2a^2/w}\le1/2$. Set
\[
 t=\sqrt{\log(m/D)},\qquad b_0=a+\sqrt w\,t,\qquad c_0=\frac{e^{-3/2}}{\sqrt{2\pi}}.
\]
Then
\begin{align}\label{eq:zero-auxiliary}
 &\E\!\left[\frac1m\sum_{i=1}^m\mathbf1_{\{|G(x_i)|\le b_0/2\}}\,\middle|\,A\right]\\
 &\qquad\le \frac2{c_0}(t+1)\frac{D+1}{\sqrt{mD}}+\frac{8w}{b_0^2}+\frac{16wD}{mb_0^2}.\notag
\end{align}
The event may impose arbitrary further constraints outside $T$.
\end{lemma}
\begin{proof}
On an enlarged Gaussian space draw $H_T\sim N(0,K_T-wI)$ and $\xi\sim N(0,I_m)$ independently, set $G_T=H_T+\sqrt w\xi$, and draw the other coordinates from their original conditional law given $G_T$. This reproduces the original full law of $G$. The outside coordinates may depend on $\xi$; no independence from them is asserted. The event inclusion alone gives
\[
 \Pp(A\mid H_T)\le\prod_{i=1}^m B_{a,w}(H_T(i)),\quad
 B_{a,w}(h)=\Phi((h-a)/\sqrt w)+\Phi((-h-a)/\sqrt w).
\]
The even function $B_{a,w}$ increases with $|h|$. If $q_H=m^{-1}\#\{i:|H_T(i)|\le b_0\}$ and $c_b=-\log B_{a,w}(b_0)$, then
\[
 \Pp(q_H\ge q\mid A)\le\min\{1,e^{D-mc_bq}\},\qquad
 \E[q_H\mid A]\le\frac{D+1}{mc_b}.
\]
Now $1-B_{a,w}(b_0)=\overline\Phi(t)-\overline\Phi(t+2a/\sqrt w)\ge\overline\Phi(t)/2$. Indeed, for $t,c\ge0$, changing variables in the Gaussian integral gives
\[
 \overline\Phi(t+c)=\int_t^\infty\phi(s+c)\,ds
 \le e^{-ct-c^2/2}\overline\Phi(t).
\]
Apply this with $c=2a/\sqrt w$ and the hypothesis
$e^{-2a^2/w}\le1/2$. Integrating the normal density over $[t,t+1/(t+1)]$ gives $\overline\Phi(t)\ge c_0 e^{-t^2/2}/(t+1)$. Since $-\log B\ge1-B$, these estimates give the first term on the right of \eqref{eq:zero-auxiliary} as an upper bound for $\E[q_H\mid A]$.

The enlarged posterior has relative entropy exactly $-\log\Pp(A)\le D$ from its prior. Data processing gives the same upper bound for the relative entropy of its $\xi$ marginal from $N(0,I_m)$. The entropy variational inequality and the Gaussian moment generating function imply
\[
 \tfrac14\E[\|\xi\|^2\mid A]\le D+\log\E e^{\|\xi\|^2/4}=D+\tfrac m2\log2,
\]
so $\E[\|\xi\|^2\mid A]\le4D+2m$. If $|G_T(i)|\le b_0/2$ and $|H_T(i)|>b_0$, then $\sqrt w|\xi_i|>b_0/2$. The averaged Markov inequality completes the proof.
\end{proof}


\begin{lemma}[Posterior output amplitude]\label{lem:zero-amplitude}
There are constants $b,C>0$, depending only on $L,a$, such that every
event $A$ implying $|G(x)|>a$ for all inputs and satisfying
$-\log\mathbb P(A)\le e^{.045d}$ obeys
\begin{equation}\label{eq:zero-amplitude}
 \mathbb E\!\left[\mathbb P_x\{|G(x)|\le b d^{(1-\beta)/2}\}
                    \,\middle|\,A\right]\le C/d
\end{equation}
for sufficiently large $d$.
\end{lemma}
\begin{proof}
Increase $-\log\mathbb P(A)$ to a deterministic upper bound $D_0\ge1$
with $D_0\le e^{.045d}$. Independently draw an auxiliary uniform sample
$T$ of size $m=\lfloor e^{.05d}\rfloor$. Let $u$ be odd with
$u/d\to1/50$, and put $w=\tau_u/2\asymp d^{-\beta}$.
Lemma~\ref{lem:gram-general} and the spectral estimates imply, with
exponentially high probability over $T$,
\[
 Q_{u,TT}\succeq (\tau_u/2)I_m,
 \qquad K_T\succeq wI_m.
\]
This event excludes repetitions. Indeed its error terms decay
exponentially because
$ .05<h(1/50)$ and $.10<\log2-h(1/4)$; the far-pair error also decays
exponentially since $u$ is a fixed positive fraction of $d$.

Lemma~\ref{lem:zero-auxiliary} applies for every such $T$, to the same
full event $A$. Here $\log(m/D_0)\ge .004d$ for all large $d$, so
$b_0^2\ge c d^{1-\beta}$. Its first error term is exponentially small;
the other terms are $O(1/d)$. Fix $b>0$ small enough that
$b d^{(1-\beta)/2}\le b_0/2$.

Although the enlarged Gaussian splitting depends on $T$, its marginal
law of the full $G$ is always the original Gaussian law, and its
conditional marginal is always the same actual posterior $G\mid A$.
Averaging the lemma over the independent auxiliary sample therefore
changes its left side exactly into the left side of
\eqref{eq:zero-amplitude}. Bad auxiliary samples contribute at most
their exponentially small probability. This proves the claim.
\end{proof}

On the sample event in Lemma~\ref{lem:zero-denominator}, both its
adaptive and rough bounds are at most $e^{.045d}$ for large $d$ when
$n\le e^{.04d}$. For the adaptive bound this uses
$D_r\le\exp[(h(1/200)+o(1))d]$ and $h(1/200)<.04$.
Consequently Lemma~\ref{lem:zero-amplitude} applies to the exact
margin-and-label posterior, uniformly over all labelings on that event.

\section{Positive subcritical exponents}
\label{sec:positive}
We now prove the lower bound for $0<\gamma<1$. The argument applies
simultaneously to every energy-minimizing sign function, so no rule
for choosing among ties is required.

\begin{theorem}[All minimum-energy signs have large error]
\label{thm:positive-subcritical}
For fixed $L\ge2$ and $a>0$, there exist $c>0$ and $d_0$ such that,
for every $d\ge d_0$ and every specified $0\le n\le e^{cd}$, with
probability at least $3/4$ over the training inputs, every minimizer of
\[
 \min\left\{\frac12\|z\|_K^2:
      |z(x)|\ge a\ \text{for every }x,\quad
      f(X_i)z(X_i)\ge a\ \text{for every }i\right\}
\]
has parity error at least $3/8$. On this event, for every fixed
$0<\gamma<1$, the actual network posterior satisfies
$\liminf_{N\to\infty}F_N(S)\ge1/16$.
\end{theorem}
\begin{proof}
Apply Lemma~\ref{lem:zero-bump} with its fixed cutoff $r=1$ and
$A=a$. Since $\tau_1\asymp1$ and $\widetilde K\preceq K$, its
interpolant has $K$-norm squared at most $C_{L,a}(d+nd^\beta)$.
By optimality this gives, except on a set of sample probability
$e^{-\Omega_L(d)}$, the simultaneous
bound
\begin{equation}\label{eq:positive-energy}
 B:=\|z\|_K^2\le C_{L,a}(d+nd^\beta)
\end{equation}
for all minimizers. The statement includes $n=0$, using the constant
comparator. If $n\le\exp[(\log(d+1))^3]$, Lemma~\ref{lem:small-cutoff}
therefore gives error $1/2-o(1)$, which is at least $3/8$ for large $d$.

For the remaining sample sizes, fix $\delta=1/20$, choose $A$ at least
as large as the constant required by
Lemma~\ref{lem:parity-squared-balance}, and use the cutoffs and notation
of Lemma~\ref{lem:adaptive-cutoff}. Work on the intersection of the
comparator event, the squared-coercivity event
\eqref{eq:wrong-square-coercivity}, and the tail-Gram event
\eqref{eq:adaptive-gram}. Its probability is at least
$1-2\delta-e^{-\Omega_L(d)}>3/4$ for large $d$.
The range of $n$ absorbs the additive $d$ in
\eqref{eq:positive-energy}, giving $B\le C_{L,a}nd^\beta$.

For an arbitrary minimizer write the orthogonal Fourier decomposition
\[
 z=P+b+t,\qquad P=z_{\le r},\quad
 b=z_{(r,u]},\quad t=z_{>u}.
\]
At each observed input where $fP<0$, training consistency implies
$(fP)_-\le|z-P|$. The squared sampling bound, the RKHS evaluation
bound, and the tail-Gram event therefore give
\[
 \frac14\|P\|_2^2
 \le\frac1n\sum_i(f(X_i)P(X_i))_-^2
 \le\frac1n\sum_i(z(X_i)-P(X_i))^2
 \le\frac Bn b_r\tau_r.
\]
Consequently,
\begin{align}
 \|P\|_2^2&\le C_{L,a}\rho^{-\beta}b_r,
       \label{eq:positive-head}\\
 \|b\|_2^2&\le B\ell_r
       \le C_{L,a}A\rho^{-\beta}H_r,
       \label{eq:positive-band}\\
 \|t\|_2^2&\le B\ell_u
       \le C_{L,a}d^\beta d^{-10}\tau_r=o(1).
       \label{eq:positive-tail}
\end{align}
All bounds are uniform over the minimizers and the specified sample
range. The middle-band norm need not vanish; it is included in the
polynomial to which the clipped-parity estimate is applied.

Put $h=P+b$. Its degree is at most $u$, and orthogonality gives
\begin{equation}\label{eq:positive-correlation-scale}
 \left(\frac ud\right)^2\|h\|_2^2
 \le C_{L,a,\delta}\left[
 \rho^{2-\beta}
 +\sqrt A\,\rho^{3/2-\beta}\sqrt{H_r}
 +A\rho^{2-\beta}H_r\right]+o(1).
\end{equation}
In particular the $\rho^{-1/2}$ loss in the tail-Gram bound has been
retained. Because $0<\beta\le1/2$, every power of $\rho$ on the
right is positive. Choose the upper cutoff fraction $\alpha$ sufficiently
small, after fixing these constants, that
$u\|h\|_2/(ad)\le1/8$ for all sufficiently large $d$.
Decrease $c$ in Lemma~\ref{lem:adaptive-cutoff} accordingly.
By \eqref{eq:positive-tail}, also $\|t\|_2/a\le1/8$.
Since $|z|\ge a$ everywhere, Lemma~\ref{lem:clipped-parity} now yields
\[
 |\mathbb E f\operatorname{sign}z|\le\frac14,
 \qquad
 \mathbb P_x\{\operatorname{sign}z(x)\ne f(x)\}\ge\frac38.
\]
This proves the assertion about every minimizer.

For any probability mixture of these minimizing signs, let $q(x)$
be the probability of the correct label and $R=\mathbb E_x(1-q(x))$.
The result just proved gives $R\ge3/8$, independently of the mixture
weights. With $F=\mathbb P_x\{q(x)\le2/3\}$, one always has
\[
 R\le\frac{1-F}{3}+F,
 \qquad F\ge\frac{3R-1}{2}\ge\frac1{16}.
\]
Finally, Proposition~\ref{prop:rate} shows that for each fixed
$0<\gamma<1$, the actual finite-width posterior puts mass $1-o(1)$
on these signs. Its draw error is therefore at least $3/8-o(1)$,
without requiring convergence of the relative weights among ties.
Apply the same elementary inequality at every $N$ and take the
lower limit to obtain $\liminf_{N\to\infty}F_N(S)\ge1/16$.
\end{proof}

\section{The fixed-margin Gaussian endpoint}
\label{sec:zero-endpoint}
We now prove the confidence lower bound for the exact $\gamma=0$
posterior. The full event $A_S$ from Section~\ref{sec:zero-conditioning}
uses the parity labels $y_i=f(X_i)$. Write $\mathbb E_S^*$ for expectation
under $G\mid A_S$, and let $\mathcal D=\mathcal D(S)$.
Its relative entropy from $N(0,K)$ is exactly $\mathcal D$.

\begin{lemma}[Gaussian quadratic entropy]\label{lem:zero-quadratic}
Suppose $Z$ is a linear image of a centered Gaussian vector under a
reference law, with covariance $T\succeq0$, trace $t$, and largest
eigenvalue $\ell>0$. Under any law of relative entropy at most $D_0$
from that reference law,
\[
 \mathbb E\|Z\|^2\le2t+4\ell D_0.
\]
\end{lemma}
\begin{proof}
The entropy variational inequality states that
$\mathbb E_Q V\le\operatorname{KL}(Q\Vert P)+\log\mathbb E_Pe^V$.
It follows immediately from nonnegativity of the relative entropy
from $e^VdP/\mathbb E_Pe^V$, first for bounded $V$ and then by
truncation. At $V=\|Z\|^2/(4\ell)$, the Gaussian moment generating
function gives
\[
 \log\mathbb E_Pe^V
 =-\frac12\sum_j\log\left(1-\frac{\lambda_j(T)}{2\ell}\right)
 \le\frac{t}{2\ell},
\]
using $-\log(1-x)\le2x$ for $0\le x\le1/2$. The claim follows.
The zero-covariance case is immediate.
\end{proof}

For any odd cutoff $r$, apply this lemma to the high-degree projection
$G_{>r}$ as a vector scaled by $M^{-1/2}$, and to its restrictions at
the $n$ training locations scaled by $n^{-1/2}$. The resulting exact
bounds are
\begin{align}
 \mathbb E_S^*\|G_{>r}\|_2^2
 &\le2\tau_r+4\lambda_{r+1}\mathcal D,
 \label{eq:zero-tail-population}\\
 \mathbb E_S^*\mathbb P_nG_{>r}^2
 &\le2\tau_r+\frac{4\|Q_{r,SS}\|\mathcal D}{n}.
 \label{eq:zero-tail-empirical}
\end{align}
Here $\|g\|_2^2=\mathbb E_xg(x)^2$ and
$\mathbb P_ng=n^{-1}\sum_i g(X_i)$. The normalized empirical covariance
has trace exactly $\tau_r$, even if locations repeat. Neither formula
asserts independence under the posterior, or log-concavity of the
globally conditioned law.

\begin{theorem}[Endpoint confidence lower bound]\label{thm:zero-endpoint}
For fixed $L\ge2$ and $a>0$, there are $c>0$ and $d_0$ such that, for
every $d\ge d_0$ and each integer $0\le n\le e^{cd}$, with probability
at least $3/4$ over independent uniform parity training inputs,
\[
 R_S:=\mathbb E_S^*\mathbb P_x\{f(x)G(x)<0\}\ge\frac38.
\]
Consequently the fraction of inputs whose correct-label posterior
probability is at most $2/3$ is at least $1/16$.
\end{theorem}
\begin{proof}
We first consider $1\le n\le\exp[(\log(d+1))^3]$.
On the exponentially high-probability event of
Lemma~\ref{lem:zero-denominator},
$\mathcal D\le C(d^2+nd^{1+\beta})$.
Let $s$ be odd with $s\sim d^{1/3}$ and put $p=G_{\le s}$.
The spectral estimates give $\tau_s=o(1)$ and
\[
 \mathcal D\lambda_{s+1}
 \le \frac{\operatorname{poly}(d)\exp[(\log(d+1))^3]}
                {\binom{d+1}{s+2}}=o(1),
\]
uniformly in this range of sample sizes. Thus
\eqref{eq:zero-tail-population} implies
$\mathbb E_S^*\|G-p\|_2^2=o(1)$.
Since $|G(x)|>a$ at every input, the signs of $G$ and $p$ can disagree
only where $|G-p|\ge a$, using any fixed sign convention at $p=0$.
Their posterior-averaged disagreement fraction is therefore $o(1)$.
Lemma~\ref{lem:ptf-parity} gives parity error $1/2-o(1)$ for every
degree-$s$ polynomial threshold, so $R_S\ge1/2-o(1)>3/8$.

For the remaining range, fix $\delta=1/20$. Choose a fixed sufficiently
large sampling constant $A$ and use the maximal odd cutoff of
Lemma~\ref{lem:adaptive-cutoff}:
\[
 A D_r\log(2D_r)\le n<A D_{r+2}\log(2D_{r+2}),
 \qquad \tfrac14(\log d)^2\le r\le\alpha d.
\]
Here $\alpha<1/256$ is a small positive constant to be chosen below;
the exponential constant $c$ is then chosen small enough to enforce
this upper bound. Write
\[
 \rho=r/d,\qquad H_\rho=\log(e/\rho),\qquad
 \tau=\tau_r,\quad\ell=\lambda_{r+1},\qquad
 B_r=C_\delta\bigl(1+\sqrt{A H_\rho/\rho}+A H_\rho\bigr).
\]
Using $r/d$ instead of $r/(d+1)$ in the adaptive lemma changes only
fixed constants.
With probability at least $1-2\delta-o(1)$, the denominator,
one-sided squared-loss sampling, and tail Gram events all hold:
\begin{align}
 \mathcal D/n&\le C d\rho^{-\beta},
 \label{eq:zero-normalized-cost}\\
 \mathbb P_n(fP)_-^2&\ge\tfrac14\|P\|_2^2
       \quad\hbox{for every degree-$r$ polynomial $P$},
 \label{eq:zero-coercivity}\\
 \|Q_{r,SS}\|&\le B_r\tau.
 \label{eq:zero-gram}
\end{align}
For the first inequality use \eqref{eq:zero-adaptive-cost},
$n\ge A D_r\log(2D_r)$, and $\tau\asymp r^{-\beta}$.
Indeed $d^2/n=o(1)$ here and
$D_r\log(r+1)/n\le C$, while the main term is
$d^{1+\beta}\tau\le C d\rho^{-\beta}$.
The other two inequalities follow from
\eqref{eq:wrong-square-coercivity} and
Lemma~\ref{lem:adaptive-cutoff}, respectively. Fix a sample on these
events for the rest of the argument.

Take $u=r+2\lceil10\log(d+1)\rceil$ and decompose the actual posterior
output into orthogonal Walsh bands,
\[
 G=p+b+t,\qquad p=G_{\le r},\quad b=G_{(r,u]},\quad t=G_{>u}.
\]
Every posterior draw has $f(X_i)G(X_i)>0$. Therefore, pointwise at
the training inputs, $(fp)_-\le|G_{>r}|$.
Equations~\eqref{eq:zero-coercivity}--\eqref{eq:zero-gram} and
\eqref{eq:zero-tail-empirical} give
\begin{equation}\label{eq:zero-head}
 \mathbb E_S^*\|p\|_2^2
 \le8\tau+16 B_r\tau\mathcal D/n.
\end{equation}
The middle band has prior trace at most $\tau$ and largest normalized
eigenvalue at most $\ell$. Lemma~\ref{lem:zero-quadratic} gives
\[
 \mathbb E_S^*\|b\|_2^2\le2\tau+4\ell\mathcal D.
\]
Orthogonality, together with the shell estimates
$n\ell/\tau\le CAH_\rho$ and
$n\lambda_{u+1}/\tau\le d^{-10}$ from
Lemma~\ref{lem:adaptive-cutoff}, now implies
\begin{align}
 \mathbb E_S^*\|p+b\|_2^2
 &\le10\tau+C\frac{\mathcal D\tau}{n}(B_r+AH_\rho),
 \label{eq:zero-low-norm}\\
 \mathbb E_S^*\|t\|_2^2
 &\le2\tau_u+4d^{-10}\frac{\mathcal D\tau}{n}.
 \label{eq:zero-final-tail}
\end{align}
No additional empirical sampling event at $u$ is needed.

Set $A_d=b_0d^{(1-\beta)/2}$, with the fixed constant $b_0>0$ of
Lemma~\ref{lem:zero-amplitude}. That lemma applies because
$\mathcal D\le e^{.045d}$ on the denominator event.
Its conclusion and \eqref{eq:zero-final-tail} yield, uniformly over
the sample sizes in question,
\begin{align}
 \mathbb E_S^*\mathbb P_x\{|G(x)|\le A_d\}&\le C/d,
 \label{eq:zero-large-outputs}\\
 \mathbb E_S^*\|t\|_2^2/A_d^2
 &\le C\{(d\rho^\beta)^{-1}+d^{-10}\rho^{-2\beta}\}=o(1).
 \label{eq:zero-tail-relative}
\end{align}
The last assertion uses $\rho\ge(\log d)^2/(4d)$ and $\beta\le1/2$.

Let $\operatorname{clip}(z)=\max(-1,\min(1,z))$.
The exact clipped-parity identity in Lemma~\ref{lem:clipped-parity},
applied to the degree-$u$ polynomial $p+b$, gives
\[
 \left|\mathbb E_x f(x)\operatorname{clip}((p+b)(x)/A_d)\right|
 \le\frac{u}{dA_d}\|p+b\|_2.
\]
Clipping is $1$-Lipschitz, and
$|\sign G-\operatorname{clip}(G/A_d)|\le\mathbf1_{\{|G|\le A_d\}}$.
Consequently Cauchy--Schwarz and the preceding bounds show
\begin{align}
 \left|\mathbb E_S^*\mathbb E_x f(x)\sign G(x)\right|
 &\le\frac{u}{dA_d}\sqrt{\mathbb E_S^*\|p+b\|_2^2}
       +\frac{\sqrt{\mathbb E_S^*\|t\|_2^2}}{A_d}+C/d.
 \label{eq:zero-correlation}
\end{align}
The second and third terms tend to zero uniformly. Since $u/r=1+o(1)$,
\eqref{eq:zero-normalized-cost} and \eqref{eq:zero-low-norm} bound
the square of the first term by
\begin{equation}\label{eq:zero-rho-powers}
 C\rho^{2-2\beta}
       \bigl(1+\rho^{-1/2}\sqrt{H_\rho}+H_\rho\bigr)+o(1).
\end{equation}
All constants here are independent of $d,n,\rho$; they may depend on
the fixed $L,a,A,\delta$.
The worst power is $\rho^{3/2-2\beta}\sqrt{H_\rho}$, which tends to
zero as $\rho\downarrow0$ for every $\beta\le1/2$.
Choose $\alpha$ small enough and then $d$ large enough that the right
side of \eqref{eq:zero-correlation} is at most $1/4$.
Thus
\[
 R_S=\frac{1-\mathbb E_S^*\mathbb E_x f(x)\sign G(x)}2\ge\frac38.
\]
The good sample event has probability at least
$1-2\delta-o(1)>3/4$. At $n=0$, global sign symmetry gives correct-label
probability $1/2$ at every input and the same risk conclusion directly.

Finally put $q_S(x)=\mathbb P\{f(x)G(x)>0\mid A_S\}$ and
$F_S=\mathbb P_x\{q_S(x)\le2/3\}$. Then
\[
 R_S=\mathbb E_x(1-q_S(x))
 \le\frac{1-F_S}{3}+F_S,
 \qquad F_S\ge\frac{3R_S-1}{2}\ge\frac1{16}.
\]
This conclusion concerns the actual posterior probability, without
identifying it with the sign of a posterior mean.
\end{proof}

\section{The Gaussian sign posterior: negative exponents}
\label{sec:negative-gamma}

Let $G\sim\mathcal N(0,K)$ on the finite cube, and condition on
\[
 \mathcal A_S=\{f(X_i)G(X_i)>0:\ 1\le i\le n\}.
\]
Write $m=\mathbb E[G\mid\mathcal A_S]$ and
$q_S(x)=\mathbb P\{f(x)G(x)>0\mid\mathcal A_S\}$.
We use the buffered threshold
\[
 \theta=\frac{13}{20},\qquad
 c_\theta=\sqrt{2\pi}\left(\frac1e-\frac7{20}\right)>0,
 \qquad F_\theta(S)=\mathbb P_x\{q_S(x)\le\theta\}.
\]
The buffer $\theta<2/3$ will preserve the strict confidence convention
when taking the finite-width limit.

\begin{lemma}[Gaussian mean and confidence]
\label{lem:negative-mean-confidence}
If $s$ is the number of distinct training inputs, then
\[
 \|m\|_K^2\le s\le n,\qquad f(X_i)m(X_i)>0.
\]
Let $\sigma_x^2=\operatorname{Var}(G(x)\mid G_S)$, with $S$ listing
the distinct inputs. If $\sigma_x>0$ and $q_S(x)>\theta$, then
$f(x)m(x)>c_\theta\sigma_x$.
For every Fourier cutoff $v$, moreover,
\begin{equation}\label{eq:negative-variance-tail}
 \mathbb P_x\{\sigma_x^2<\tau_v/2\}
       \le\frac{2n\ell_v}{\tau_v},
 \qquad
 \|m-\Pi_{\le v}m\|_2^2\le n\ell_v.
\end{equation}
\end{lemma}

\begin{proof}
Whiten the training vector as $G_S=K_{SS}^{1/2}\xi$.
The sign constraints form a homogeneous cone in $\xi$. Since a standard
Gaussian's radius is independent of its direction, conditioning leaves
$\mathbb E\|\xi\|_2^2=s$. Jensen's inequality and Gaussian regression
therefore give
\[
 m_S^{\mathsf T}K_{SS}^{-1}m_S\le s,
 \qquad m=K_{\mathcal X,S}K_{SS}^{-1}m_S.
\]
The squared RKHS norm of the regression vector is the first quadratic
form, and its signed training values are positive.

The conditional scalar variable $Y=f(x)G(x)$ has a log-concave density:
it is a marginal of a Gaussian restricted to a convex set, by
Pr\'ekopa's marginal theorem \cite{prekopa1973}. Its cdf $H$ is also
log-concave. For $\mu=\mathbb E Y$, Jensen's inequality and the probability
integral transform give
\[
 \log H(\mu)\ge\mathbb E\log H(Y)=-1,
 \qquad \mathbb P\{Y>\mu\}\le1-e^{-1}.
\]
Conditioning first on $G_S$ expresses the density of $Y$ as a mixture
of Gaussian densities of common variance $\sigma_x^2$. It is bounded
above by $(\sqrt{2\pi}\sigma_x)^{-1}$. If $q_S(x)>\theta$, this first
forces $\mu>0$, and then
\[
 \theta<q_S(x)\le1-e^{-1}
                  +\frac{\mu}{\sqrt{2\pi}\sigma_x}.
\]
This proves the confidence margin.

To bound residual variances, write $G=G_{\le v}+G_{>v}$ independently.
Conditioning additionally on the entire low-degree field can only
reduce a Gaussian conditional variance. The remaining covariance is
\[
 Q_v-Q_{v,\mathcal X,S}Q_{v,SS}^{\dagger}Q_{v,S,\mathcal X}.
\]
The subtracted matrix is positive semidefinite, bounded by $Q_v$, and
has rank at most $s$. Its normalized trace is at most $n\ell_v$.
Markov's inequality proves the first bound in
\eqref{eq:negative-variance-tail}; this automatically includes the
observed points, where the full residual variance vanishes.
The second bound follows from Fourier orthogonality and
$\|m\|_K^2\le n$.
\end{proof}

\begin{theorem}[An exponential lower bound for the sign posterior]
\label{thm:negative-all-n}
For every fixed depth $L\ge2$, there is $c_L>0$ such that, for every
sufficiently large $d$ and each integer $0\le n\le e^{c_Ld}$,
\[
 \mathbb P_S\{F_{13/20}(S)\ge3/8\}\ge3/4.
\]
Consequently the same lower bound holds for the $\liminf$ of the
finite-width strict-$2/3$ confidence-failure fraction, for each fixed
$\gamma<0$. The width limit is taken at fixed $d$ and training sample.
\end{theorem}

\begin{proof}
We give the cutoff estimates explicitly. Fix the sampling-failure
probability at a sufficiently small universal constant and let $A$ be
the resulting constant in Lemma~\ref{lem:parity-squared-balance}.
Choose the largest odd $r$ such that
\[
 n\ge A D_r\log(2D_r).
\]
First suppose $n\ge\exp((\log(d+1))^3)$. Then the adaptive cutoff
estimates of Lemma~\ref{lem:adaptive-cutoff} give
$r\ge(\log d)^2/4$ for large $d$. By reducing $c_L$ we may ensure
$r\le\alpha d$, for a fixed $\alpha<1/256$ to be chosen below.
Write $\rho=r/d$ and $L_\rho=\log(e/\rho)$.
Maximality and the sharp paired eigenvalue estimates give
\begin{equation}\label{eq:negative-cutoff-bounds}
 \frac{n\ell_r}{\tau_r}\le C_A L_\rho,
 \qquad
 \|Q_{r,SS}\|_{\mathrm{op}}\le b_r\tau_r,
 \qquad b_r=C_{L,A}\left(1+\sqrt{L_\rho/\rho}+L_\rho\right),
\end{equation}
with probability at least $1-\delta$, where $\delta$ is a sufficiently
small fixed constant. Indeed,
$n<A D_{r+2}\log(2D_{r+2})$,
$D_{r+2}\asymp\binom{d+1}{r+2}$, and
$\log D_{r+2}\le C r\log(ed/r)$. The ratio estimate
\eqref{eq:ratio} proves the first bound. For the Gram bound use
Lemma~\ref{lem:gram-general}, with
\[
 \frac{nd\,s_{2,r}}{\tau_r^2}\le C_A L_\rho/\rho,
 \qquad \frac{n\ell_r}{\tau_r}\le C_A L_\rho.
\]
The far-pair and collision errors vanish uniformly in this range when
$c_L$ is sufficiently small.

Let $p=\Pi_{\le r}m$. On the squared-coercivity event,
training consistency and the RKHS evaluation inequality imply
\[
 \frac14\|p\|_2^2
 \le\frac1n\sum_i(f(X_i)p(X_i))_-^2
 \le\frac1n\sum_i(m(X_i)-p(X_i))^2
 \le\frac{\|m\|_K^2}{n}\|Q_{r,SS}\|_{\mathrm{op}}
 \le b_r\tau_r.
\]
Choose an odd $u=r+O(\log d)$, with a sufficiently large fixed implicit
constant. Compare the sharp eigenvalue bounds at the two endpoints:
$\ell_u/\ell_r\le C_L\binom{d+1}{r+2}/\binom{d+1}{u+2}$,
up to a bounded power ratio. The binomial ratio is geometrically small
in $u-r$, uniformly for $r\le\alpha d$. The sharp level masses
and the bound $r\ge(\log d)^2/4$ therefore give
\begin{equation}\label{eq:negative-high-cutoff}
 \tau_u=(1-o(1))\tau_r,\qquad n\ell_u=o(\tau_r).
\end{equation}
These estimates are uniform over the sample-size range under
consideration: the removed band has trace $O((\log d)/r)\tau_r$,
and the geometric decay over $O(\log d)$ levels absorbs all polynomial
factors in \eqref{eq:negative-cutoff-bounds}.
Put $h=\Pi_{\le u}m$ and $t=m-h$. Orthogonality yields
\begin{equation}\label{eq:negative-low-norm}
 \|h\|_2^2\le4b_r\tau_r+n\ell_r
             \le C_{L,A}\tau_r(b_r+L_\rho),
 \qquad \|t\|_2^2=o(\tau_r).
\end{equation}

Lemma~\ref{lem:negative-mean-confidence} and
\eqref{eq:negative-high-cutoff} show that outside an $o(1)$ fraction
of inputs, $\sigma_x^2\ge\tau_r/3$ and
$|t(x)|\le c_\theta\sqrt{\tau_r}/(2\sqrt3)$.
Set $a_r=c_\theta\sqrt{\tau_r}/(2\sqrt3)$.
At every remaining input where $q_S(x)>\theta$, we have
$f(x)h(x)>a_r$. Consequently
\begin{equation}\label{eq:negative-confidence-clipping}
 \mathbb E f\,\operatorname{clip}(h/a_r,-1,1)
 \ge1-2F_\theta(S)-o(1).
\end{equation}
For any polynomial $h$ of degree at most $u$,
Lemma~\ref{lem:clipped-parity} gives
\[
 \left|\mathbb E f\,\operatorname{clip}(h/a_r,-1,1)\right|
 \le\frac{\|\Delta h\|_1}{a_r d}
 \le\frac{u\|h\|_2}{a_r d},
 \qquad \Delta=\sum_{j=1}^dD_j.
\]
For clarity, the first inequality follows directly by setting $q=fh$
and letting $R$ average the one-coordinate flips. Then
$Rq=-q+(2/d)f\Delta h$. Stationarity, Jensen's inequality and the
triangle inequality imply
$|\mathbb E|q+a_r|-\mathbb E|q-a_r||
\le(2/d)\|\Delta h\|_1$.
Use $\operatorname{clip}(q/a_r,-1,1)
=(|q+a_r|-|q-a_r|)/(2a_r)$.
The second inequality uses the Walsh multiplier bound
$\|\Delta h\|_2\le u\|h\|_2$.
Thus \eqref{eq:negative-low-norm} bounds the absolute value in
\eqref{eq:negative-confidence-clipping} by
\[
 C_{L,A,c_\theta}\rho
       \sqrt{1+\sqrt{L_\rho/\rho}+L_\rho}+o(1),
\]
which tends to zero as $\rho\downarrow0$. Choose $\alpha$ small enough
that it is at most $1/8$ for $0<\rho\le\alpha$, and then take $d$
large enough to absorb the vanishing errors. This proves
$F_\theta(S)\ge3/8$. The sampling and Gram events have combined
probability at least $3/4$ after the initial fixed choices.

It remains to cover smaller samples, including cases where the cutoff
$r$ does not exist. If $n<\exp((\log(d+1))^3)$, choose an odd
$v\asymp d^{1/3}$. The sharp spectrum gives
$n\ell_v/\tau_v=o(1)$, because
$\log\binom{d+1}{v+2}\asymp v\log(d/v)$ dominates $\log(n+1)$.
Let $p=\Pi_{\le v}m$. By
Lemma~\ref{lem:negative-mean-confidence}, outside an $o(1)$ fraction
of inputs, the residual variance is at least $\tau_v/2$ and
$|m-p|<c_\theta\sqrt{\tau_v}/(2\sqrt2)$.
Therefore confidence above $\theta$ implies $fp>0$ there.
Lemma~\ref{lem:ptf-parity}, with $v=o(\sqrt d)$, gives
$F_\theta(S)\ge1/2-o(1)\ge3/8$ for every such training sample.
For $n=0$ the conclusion also follows immediately from $q_S=1/2$.
Finally, the conditional width limit and $\theta<2/3$ give the asserted
finite-width $\liminf$ statement.
\end{proof}


\section{Proof of the main theorem}
The positive-exponent, zero-exponent, and negative-exponent conclusions
proved above each hold with probability at least $3/4$ for every specified
$n\le e^{cd}$. Decrease $c$ and increase $d_0$ to satisfy their finitely
many requirements. The width-limit propositions transfer these conclusions
to the stated neural-network posterior.

For clarity, when the posterior's mean classification error is $R_N$,
the confidence fraction always satisfies
\[
 R_N=\mathbb E_x(1-q_{N,S}(x))
 \le\frac{1-F_N}{3}+F_N,
 \qquad F_N\ge\frac{3R_N-1}{2}.
\]
The nonnegative-exponent arguments give $\liminf R_N\ge3/8$, hence the
required $1/16$. The negative-exponent argument instead directly bounds
the fraction with limiting confidence at most $13/20<2/3$, so convergence
of each of the finitely many posterior label probabilities gives the same
conclusion. This proves Theorem~\ref{thm:main}.

\appendix
\section{A small polynomial-weight parity network}\label{sec:paritycircuit}
For completeness, full parity on $k$ bits has a one-hidden-layer
$\tanh$ realization of width $k+1$, with weights $O(k\log(k+1))$
and output margin greater than $1/2$. Let
\[
 t(x)=\frac{k+\sum_i x_i}{2}\in\{0,\ldots,k\},\quad
 y_j=(-1)^{k-j},\quad A=\log(16(k+1)).
\]
The step functions $s_j(t)=\sign(t-j+1/2)$ satisfy
\[
 y_t=c+\sum_{j=1}^k y_js_j(t),\qquad c=(y_0+y_k)/2\in\{0,1\}.
\]
Replace each step by $h_j=\tanh(A(t-j+1/2))$. Since its argument
has absolute value at least $A/2$, its error is at most $2e^{-A}$.
A constant neuron $h_0=\tanh1$ supplies the term $c$ using output
weight $c/\tanh1$. The resulting output differs from parity by
less than $1/8$. First-layer input weights are $A/2$ and bias
weights have magnitude at most $Ak/2$; readout weights have
magnitude at most $2$.

At depth $L\ge3$, the next hidden layer sums these features and
applies $\tanh$, producing the correct sign with magnitude at least
$\tanh(7/8)>1/2$. Each additional hidden layer can preserve this
margin by a single weight $2$, since $\tanh(2u)>1/2$ whenever
$u>1/2$. A final readout weight $1$ completes the construction.
Unused neurons are set to zero. Thus the same polynomial bounds
hold at every fixed depth.
For any other fixed margin $a>0$, multiplying the final readout by $2a$
gives margin greater than $a$ with the same polynomial complexity bounds.



\begin{thebibliography}{9}
\bibitem{ABFR}
J. Aspnes, R. Beigel, M. Furst, and S. Rudich.
The expressive power of voting polynomials.
\emph{Combinatorica} 14(2):135--148, 1994.
\href{https://www.cs.yale.edu/homes/aspnes/papers/stoc91voting.pdf}{Author manuscript}.
\bibitem{flajolet2009}
P. Flajolet and R. Sedgewick. \emph{Analytic Combinatorics}.
Cambridge University Press, 2009. Chapter VI.
\bibitem{Lee2023}
J. R. Lee. Spectral hypergraph sparsification via chaining.
\href{https://arxiv.org/abs/2209.04539}{arXiv:2209.04539},
Section 2.4, Lemma 2.12 and equation (2.20).
\bibitem{tropp2008}
J. A. Tropp. Norms of random submatrices and sparse approximation.
\emph{Comptes Rendus Math\'ematique} 346:1271--1274, 2008.
\href{https://tropp.caltech.edu/papers/Tro08-Norms-Random-v3.pdf}{Author manuscript},
Theorem 1.1.
\bibitem{prekopa1973}
A. Pr\'ekopa. On logarithmic concave measures and functions.
\emph{Acta Scientiarum Mathematicarum} 34:335--343, 1973.
\end{thebibliography}
\end{document}
